% LISTA DE ABREVIATURAS E SIGLAS -----------------------------------------------
% Conforme NBR 14724 (Apresentação de trabalhos acadêmicos). Ordem alfabética
% das siglas; expansão por extenso à direita. Termos próprios (AAPL, FinBERT,
% LightGBM) omitidos por serem nomes de entidades/ticker/software.

\begin{siglas}
    \item[ADR] Registro de decisão arquitetural (\textit{Architecture Decision Record})
    \item[API] Interface de programação de aplicações (\textit{Application Programming Interface})
    \item[AR(1)] Modelo autorregressivo de ordem 1 (\textit{Autoregressive of order 1})
    \item[ARCH] Heterocedasticidade condicional autorregressiva (\textit{Autoregressive Conditional Heteroskedasticity})
    \item[CAMELS] Indicadores bancários de adequação de capital, qualidade dos ativos, gestão, resultados, liquidez e sensibilidade ao risco de mercado
    \item[CNN] Rede neural convolucional (\textit{Convolutional Neural Network})
    \item[CPU] Unidade central de processamento (\textit{Central Processing Unit})
    \item[CQR] Regressão quantílica conformalizada (\textit{Conformalized Quantile Regression})
    \item[CRPS] \textit{Continuous Ranked Probability Score}
    \item[CVaR] Valor em risco condicional (\textit{Conditional Value-at-Risk})
    \item[DM] Teste de Diebold-Mariano
    \item[EWMA] Média móvel exponencialmente ponderada (\textit{Exponentially Weighted Moving Average})
    \item[FHS] Simulação histórica filtrada (\textit{Filtered Historical Simulation})
    \item[FWER] Taxa de erro por família (\textit{Family-Wise Error Rate})
    \item[GARCH] Modelo autorregressivo de heterocedasticidade condicional generalizado (\textit{Generalized Autoregressive Conditional Heteroskedasticity})
    \item[GBM] Máquina de \textit{gradient boosting} (\textit{Gradient Boosting Machine})
    \item[GJR-GARCH] GARCH assimétrico de Glosten, Jagannathan e Runkle
    \item[GPU] Unidade de processamento gráfico (\textit{Graphics Processing Unit})
    \item[HAC] Estimador consistente sob heterocedasticidade e autocorrelação (\textit{Heteroskedasticity and Autocorrelation Consistent})
    \item[HLN] Correção Harvey-Leybourne-Newbold para o teste DM
    \item[i.i.d.] Independentes e identicamente distribuídas
    \item[LOCO] Remoção de um grupo de covariáveis com retreino (\textit{Leave-One-Covariate-Out})
    \item[LSTM] Memória de longo e curto prazo (\textit{Long Short-Term Memory})
    \item[MA] Média móvel (\textit{Moving Average})
    \item[MACD] Convergência e divergência de médias móveis (\textit{Moving Average Convergence Divergence})
    \item[MAE] Erro absoluto médio (\textit{Mean Absolute Error})
    \item[MAPE] Erro percentual absoluto médio (\textit{Mean Absolute Percentage Error})
    \item[MCS] Conjunto de modelos de confiança (\textit{Model Confidence Set})
    \item[ML] Aprendizado de máquina (\textit{Machine Learning})
    \item[MLP] Perceptron multicamadas (\textit{Multilayer Perceptron})
    \item[MPIW] Largura média do intervalo de predição (\textit{Mean Prediction Interval Width})
    \item[MQRNN] Rede recorrente quantílica multi-horizonte (\textit{Multi-Horizon Quantile Recurrent Neural Network})
    \item[NexCP] Predição conformal ponderada para dados não permutáveis (\textit{Nonexchangeable Conformal Prediction})
    \item[OHLCV] Abertura, máxima, mínima, fechamento e volume (\textit{Open, High, Low, Close, Volume})
    \item[OLS] Mínimos quadrados ordinários (\textit{Ordinary Least Squares})
    \item[p.p.] Pontos percentuais
    \item[PICP] Probabilidade de cobertura do intervalo de predição (\textit{Prediction Interval Coverage Probability})
    \item[QLIKE] Perda de quase-verossimilhança (\textit{Quasi-Likelihood})
    \item[RMSE] Raiz do erro quadrático médio (\textit{Root Mean Squared Error})
    \item[ROCm] Plataforma de computação em GPU da AMD (\textit{Radeon Open Compute})
    \item[RSI] Índice de força relativa (\textit{Relative Strength Index})
    \item[SHA-256] Algoritmo de \textit{hash} criptográfico de 256 bits (\textit{Secure Hash Algorithm 256-bit})
    \item[TFT] Temporal Fusion Transformer
    \item[TPE] Estimador de Parzen estruturado em árvore (\textit{Tree-structured Parzen Estimator})
    \item[UTC] Tempo Universal Coordenado (\textit{Coordinated Universal Time})
    \item[VaR] Valor em risco (\textit{Value-at-Risk})
    \item[VSN] Rede de seleção de variáveis (\textit{Variable Selection Network})
\end{siglas}
